% Section 02 - Installation et démarrage
\section{Installation et premier démarrage}

\subsection{Téléchargement de l'application}

\subsubsection{Où télécharger ?}

L'application est distribuée sous forme d'installateur pour chaque système d'exploitation :

\begin{itemize}[leftmargin=*]
    \item \textbf{Windows} : Fichier \texttt{FMSP-Releves-Setup-x.x.x.exe}
    \item \textbf{macOS} : Fichier \texttt{FMSP-Releves-x.x.x.dmg}
    \item \textbf{Linux} : Fichier \texttt{FMSP-Releves-x.x.x.AppImage} ou \texttt{.deb}
\end{itemize}

\astuce{Vérifiez toujours que vous téléchargez la dernière version depuis la source officielle.}

\subsection{Installation sur Windows}

\begin{enumerate}[leftmargin=*]
    \item Double-cliquez sur le fichier \texttt{FMSP-Releves-Setup-x.x.x.exe}
    \item Si Windows SmartScreen affiche un avertissement :
    \begin{itemize}
        \item Cliquez sur \bouton{Informations complémentaires}
        \item Puis sur \bouton{Exécuter quand même}
    \end{itemize}
    \item Suivez l'assistant d'installation :
    \begin{itemize}
        \item Acceptez la licence
        \item Choisissez le répertoire d'installation (défaut : \texttt{C:\textbackslash Program Files\textbackslash FMSP-Releves})
        \item Cochez "Créer un raccourci sur le bureau" si souhaité
        \item Cliquez sur \bouton{Installer}
    \end{itemize}
    \item Attendez la fin de l'installation (environ 30 secondes)
    \item Cliquez sur \bouton{Terminer} pour lancer l'application
\end{enumerate}

\begin{figure}[h]
    \centering
    % \includegraphics[width=0.7\textwidth]{images/installation_windows.png}
    \fbox{\parbox{0.7\textwidth}{\centering [CAPTURE D'ÉCRAN : ASSISTANT D'INSTALLATION WINDOWS]}}
    \caption{Assistant d'installation sous Windows}
\end{figure}

\subsection{Installation sur macOS}

\begin{enumerate}[leftmargin=*]
    \item Double-cliquez sur le fichier \texttt{FMSP-Releves-x.x.x.dmg}
    \item Une fenêtre s'ouvre montrant l'icône de l'application
    \item Glissez-déposez l'icône vers le dossier \texttt{Applications}
    \item Éjectez le volume DMG
    \item Ouvrez le dossier \texttt{Applications}
    \item Double-cliquez sur \texttt{FMSP-Releves}
    \item Si macOS affiche "Application non vérifiée" :
    \begin{itemize}
        \item Allez dans \menu{Préférences Système} > \menu{Sécurité et confidentialité}
        \item Cliquez sur \bouton{Ouvrir quand même}
    \end{itemize}
\end{enumerate}

\subsection{Installation sur Linux}

\subsubsection{Méthode 1 : AppImage (recommandé)}

\begin{enumerate}[leftmargin=*]
    \item Rendez le fichier exécutable :
    \begin{verbatim}
chmod +x FMSP-Releves-x.x.x.AppImage
    \end{verbatim}
    \item Lancez l'application :
    \begin{verbatim}
./FMSP-Releves-x.x.x.AppImage
    \end{verbatim}
\end{enumerate}

\subsubsection{Méthode 2 : Paquet Debian (.deb)}

\begin{enumerate}[leftmargin=*]
    \item Installez le paquet :
    \begin{verbatim}
sudo dpkg -i FMSP-Releves-x.x.x.deb
    \end{verbatim}
    \item Si des dépendances manquent :
    \begin{verbatim}
sudo apt-get install -f
    \end{verbatim}
    \item Lancez depuis le menu des applications ou en ligne de commande :
    \begin{verbatim}
fmsp-releves
    \end{verbatim}
\end{enumerate}

\subsection{Premier démarrage}

\subsubsection{Écran de bienvenue}

Au premier lancement, l'application affiche un écran de bienvenue vous guidant à travers les étapes initiales.

\begin{figure}[h]
    \centering
    % \includegraphics[width=0.8\textwidth]{images/ecran_bienvenue.png}
    \fbox{\parbox{0.8\textwidth}{\centering [CAPTURE D'ÉCRAN : ÉCRAN DE BIENVENUE]}}
    \caption{Écran de bienvenue au premier démarrage}
\end{figure}

\subsubsection{Configuration initiale de l'établissement}

Avant de générer vos premiers documents, configurez les informations de votre établissement :

\begin{enumerate}[leftmargin=*]
    \item Cliquez sur le menu \menu{Configurer les entêtes}
    \item Remplissez les informations suivantes :

    \begin{tcolorbox}[colback=blue!5!white,colframe=blue!75!black,title=Informations requises]
    \begin{itemize}[leftmargin=*]
        \item \champ{Nom de l'établissement} : Ex. "Faculté de Médecine et Sciences Pharmaceutiques"
        \item \champ{Adresse complète} : Rue, ville, code postal, pays
        \item \champ{Téléphone} : Numéro de contact
        \item \champ{Email} : Adresse email officielle
        \item \champ{Site web} : URL du site (optionnel)
        \item \champ{Logo} : Image PNG ou JPG (max 2 Mo, résolution recommandée : 800x200 px)
    \end{itemize}
    \end{tcolorbox}

    \item Cliquez sur \bouton{Enregistrer les paramètres}
\end{enumerate}

\begin{figure}[h]
    \centering
    % \includegraphics[width=\textwidth]{images/config_entetes.png}
    \fbox{\parbox{0.9\textwidth}{\centering [CAPTURE D'ÉCRAN : FORMULAIRE DE CONFIGURATION DES ENTÊTES]}}
    \caption{Configuration des entêtes de l'établissement}
\end{figure}

\subsection{Activation de la licence}

\subsubsection{Mode démo vs Mode activé}

\begin{table}[h]
\centering
\begin{tabular}{|l|c|c|}
\hline
\textbf{Fonctionnalité} & \textbf{Mode démo} & \textbf{Mode activé} \\
\hline
Générer les relevés & ✓ & ✓ \\
Générer les attestations & ✓ & ✓ \\
Config des relevés & ✓ & ✓ \\
Configurer les entêtes & ✓ & ✓ \\
Historique des docs & ✓ & ✓ \\
\textbf{QR Codes sur PDF} & \textbf{✗} & \textbf{✓} \\
Placement QR sur Documents & ✓ & ✓ \\
Aide & ✓ & ✓ \\
Paramètres & ✓ & ✓ \\
\hline
\end{tabular}
\caption{Comparaison mode démo vs mode activé}
\end{table}

\attention{Le menu "QR Codes sur PDF" est désactivé en mode démo. Pour l'activer, vous devez entrer une clé de licence valide.}

\subsubsection{Procédure d'activation}

\begin{enumerate}[leftmargin=*]
    \item Cliquez sur le menu \menu{Paramètres}
    \item Dans l'onglet \textbf{Licence}, vous verrez l'état actuel (Démo ou Activé)
    \item Entrez votre clé de licence dans le champ prévu
    \item Cliquez sur \bouton{Activer}
    \item Si la clé est valide, un message de confirmation s'affiche
    \item Redémarrez l'application pour appliquer l'activation
\end{enumerate}

\begin{figure}[h]
    \centering
    % \includegraphics[width=0.8\textwidth]{images/activation_licence.png}
    \fbox{\parbox{0.8\textwidth}{\centering [CAPTURE D'ÉCRAN : ÉCRAN D'ACTIVATION DE LA LICENCE]}}
    \caption{Interface d'activation de la licence}
\end{figure}

\subsection{Structure des dossiers de l'application}

Après installation, l'application crée automatiquement les dossiers suivants sur votre système :

\subsubsection{Windows}

\begin{verbatim}
C:\Users\[VotreNom]\AppData\Roaming\FMSP-Releves\
├── configs\              # Configurations de classes sauvegardées
├── themes\               # Thèmes personnalisés
├── history\              # Historique des documents générés
└── cache\                # Cache temporaire
\end{verbatim}

\subsubsection{macOS}

\begin{verbatim}
~/Library/Application Support/FMSP-Releves/
├── configs/
├── themes/
├── history/
└── cache/
\end{verbatim}

\subsubsection{Linux}

\begin{verbatim}
~/.config/FMSP-Releves/
├── configs/
├── themes/
├── history/
└── cache/
\end{verbatim}

\info{Ces dossiers contiennent vos données personnalisées. N'hésitez pas à en faire des sauvegardes régulières.}

\subsection{Vérification de l'installation}

Pour vérifier que l'application fonctionne correctement :

\begin{enumerate}[leftmargin=*]
    \item L'interface principale s'affiche correctement avec les 9 menus dans la barre latérale
    \item Le menu \menu{Aide} affiche la documentation
    \item Le menu \menu{Paramètres} montre l'état de la licence
    \item Vous pouvez accéder à tous les menus (sauf "QR Codes sur PDF" si en mode démo)
\end{enumerate}

\begin{tcolorbox}[colback=red!5!white,colframe=red!75!black,title=En cas de problème au démarrage]
Si l'application ne se lance pas correctement :
\begin{itemize}[leftmargin=*]
    \item Vérifiez que votre système répond à la configuration minimale
    \item Désinstallez et réinstallez l'application
    \item Vérifiez les logs d'erreur dans le dossier de l'application
    \item Consultez le menu \menu{Aide} > \menu{Rapporter un problème}
\end{itemize}
\end{tcolorbox}

\subsection{Désinstallation}

\subsubsection{Windows}

\begin{enumerate}[leftmargin=*]
    \item Panneau de configuration > Programmes et fonctionnalités
    \item Recherchez "FMSP-Releves"
    \item Clic droit > Désinstaller
    \item Suivez l'assistant de désinstallation
\end{enumerate}

\subsubsection{macOS}

\begin{enumerate}[leftmargin=*]
    \item Ouvrez le dossier Applications
    \item Glissez l'icône FMSP-Releves vers la Corbeille
    \item Videz la Corbeille
\end{enumerate}

\subsubsection{Linux}

Pour AppImage : supprimez simplement le fichier .AppImage

Pour paquet Debian :
\begin{verbatim}
sudo apt-get remove fmsp-releves
\end{verbatim}

\attention{La désinstallation ne supprime pas automatiquement vos données personnelles (configurations, thèmes, historique). Ces dossiers doivent être supprimés manuellement si souhaité.}

\newpage
